\documentclass[11pt]{article}

\usepackage[margin=1in]{geometry}
\usepackage[T1]{fontenc}
\usepackage[utf8]{inputenc}
\usepackage{lmodern}
\usepackage{microtype}
\usepackage{amsmath,amssymb}
\usepackage{booktabs}
\usepackage{tabularx}
\usepackage{array}
\usepackage{enumitem}
\usepackage{pifont}
\usepackage{xcolor}
\usepackage[colorlinks=true,linkcolor=blue!50!black,citecolor=blue!50!black,urlcolor=blue!50!black]{hyperref}

\setlist{itemsep=2pt,topsep=4pt}
\newcommand{\yes}{\ding{51}}
\newcommand{\no}{\ding{55}}
\newcommand{\Ltwo}{L\textsuperscript{2}}
\newcolumntype{L}{>{\raggedright\arraybackslash}X}

\title{Taking contrastive learning to time series:\\what transfers from the paper (goals 1--3)}
\author{}
\date{}

\begin{document}
\maketitle

\begin{abstract}
\noindent
This note summarises three rounds of experiments that extend Baptista, Stuart \& Tran,
\emph{A Mathematical Perspective on Contrastive Learning}~\cite{bst}, from its Gaussian,
linear-encoder setting to forecasting time series. We relax one assumption at a time:
non-Gaussian synthetic data (goal~1), nonlinear encoders (goal~2) and a real weather
record (goal~3). The paper's probabilistic reading of contrastive learning, and its
\Ltwo{} result, survive every relaxation. Its cosine-versus-\Ltwo{} contrast and its
joint-loss penalty turn out to be properties of linear encoders. Off the paper's
setting, the things that matter most are whether the score is bounded, the bank of
real futures the forecast is built from, and the amount of data.
\end{abstract}

\noindent
Detailed numbers are in the per-goal write-ups in \texttt{docs/}:
\texttt{goal1\_nonlinear\_domain\_strategy.md} (the scoring method),
\texttt{goal1\_results.md}, \texttt{goal2\_results.md} and \texttt{goal3\_results.md}.

%======================================================================
\section{The question}

The paper reads contrastive learning (CLIP and its variants) as \emph{learning a
conditional distribution}. A trained model with score $s(\text{past},\text{future})$
defines a forecast:
\begin{quote}
take a collection of real futures and reweight each one by $\exp(\text{score})$.
\end{quote}
The paper works this out exactly when everything is Gaussian and the encoders are
linear. In that setting it makes sharp predictions about which loss gets what right
(Table~\ref{tab:predictions}).

\begin{table}[h]
\centering
\small
\begin{tabularx}{\textwidth}{@{}lL@{}}
\toprule
& The paper's prediction (Gaussian data, linear encoders) \\
\midrule
P1 & The standard dot-product score (``cosine'', unnormalised) gets the \textbf{average}
     forecast right but the \textbf{spread} wrong: its uncertainty is stuck at ``knowing
     nothing'' (Cor.~5.2). \\
P2 & The paper's new \textbf{\Ltwo{} score} (negative squared distance between embeddings),
     trained in the forecasting direction, gets \textbf{both} the average and the spread
     right (Thm~5.3, Cor.~5.4). \\
P3 & For the dot product, weighting the two directions (past $\to$ future vs
     future $\to$ past) does not change the answer. \\
P4 & The \textbf{joint} loss (one softmax over all pairs) shrinks the forecast towards the
     average, so it is a worse conditional forecaster than the one-way loss, but better at
     the marginals (Thm~5.6, Cor.~5.7). \\
P5 & Retrieval picks the most likely candidate under the model (Thm~4.2), so how well it
     can work is capped by how much the past says about the future. \\
\bottomrule
\end{tabularx}
\caption{The predictions we tested.}
\label{tab:predictions}
\end{table}

We asked how far these carry over once the setting stops being the paper's,
relaxing one assumption at a time.
\begin{enumerate}
\item \textbf{Goal 1: non-Gaussian data, still linear encoders.} Three synthetic series,
  each with a dial whose zero point is the paper's Gaussian case:
  \begin{itemize}
  \item heavy tails (Student-$t$; occasional wild values);
  \item a bent Gaussian process (skewed futures, nonlinear relation between past and future);
  \item a double-well process (the series jumps between two levels, so the future can
        have two separate likely outcomes).
  \end{itemize}
  Each has an exact ``perfect forecaster'' to compare against.
\item \textbf{Goal 2: nonlinear (MLP) encoders}, on the same three series.
\item \textbf{Goal 3: real data.} Jena weather station, 2020, 10-minute readings:
  \begin{itemize}
  \item forecast the next 6\,h of temperature from the last day of six variable groups;
  \item a cross-modal setup, closer to how CLIP is used: future rain from past pressure,
        wind and moisture.
  \end{itemize}
  There is no perfect forecaster here.
\end{enumerate}
Every round ran the full grid \{dot product, \Ltwo\} $\times$ \{one-way, two-way, joint\},
plus normalised cosine (real CLIP) from goal~2 on. The comparison points were regression
models: MSE plus a noise band, and from goal~2 on a Gaussian model with a past-dependent
band width (``Gaussian NLL'').

%======================================================================
\section{How we scored forecasts}

The earlier test bed scored models by comparing internal matrices against the paper's
closed forms, or by a cosine similarity between predicted and true futures. Neither works
once there is no closed form, and the similarity score ignores uncertainty, which is the
paper's main subject. We replaced both with a \textbf{black-box} procedure that works for
any model.

\begin{itemize}
\item \textbf{Turn the model into a forecast} exactly as the paper defines it (its eq.~31):
  keep a ``bank'' of real futures and resample them with weights $\exp(\text{score})$.
  This needs nothing but score calls: no closed form, no knowledge of the encoder.
\item \textbf{Score the forecast} with the energy score~\cite{gneiting}, a standard measure
  for probabilistic forecasts. It rewards being close to what happened and being honest
  about uncertainty; a point forecast is simply a forecast with no spread. Alongside it we
  report band coverage (does the 90\% band catch the truth 90\% of the time?), per-horizon
  errors, and how many bank futures a forecast actually uses (a collapse check).
\end{itemize}

On the Gaussian test bed this procedure reproduces the paper's closed-form answers to
within sampling error, and P1 shows up in it directly: cosine has the perfect forecaster's
average error but 98\% coverage at nominal 90\%. On real data the procedure needs only the
observed future, so the same yardstick runs from the paper's setting all the way to the
weather station. We think this scoring setup is a methodological contribution in its own
right.

%======================================================================
\section{What we found, round by round}

\subsection{Goal 1: non-Gaussian data, linear encoders}

\begin{itemize}
\item \textbf{The paper's Gaussian results hold almost exactly where they apply.}
  Linear \Ltwo{} sat on the theoretical optimum, and its retrieval matched the best
  achievable. The old evaluation had reported it about $5\times$ too low, because it scored
  retrieval with a similarity the model was never trained on. Cosine's apparent ``gap to
  theory'' was mostly slow drift in directions that are pure noise; measured in the units
  the forecast depends on, it was within about 5\% of the theory. Retrieval is low for
  every model because the past says little about the future in this setup (P5).
\item \textbf{\Ltwo{} stays within 0--5\% of the perfect forecaster on all three kinds of
  non-Gaussian data}, with linear encoders. A Gaussian regression model lands at +11--15\%
  on the bent and double-well series. The reason is that a contrastive forecast never
  invents futures, it reweights real ones. Skew and two separate outcomes therefore come
  from the data, and the encoder only has to choose the weights.
\item \textbf{Cosine's too-wide bands (P1) get worse off-Gaussian, and on heavy tails it
  breaks completely.} At the heaviest tails the forecast collapses onto a handful of
  extreme training futures and is about $10\times$ worse than perfect.
\item \textbf{One-way vs two-way made no difference (P3)}, for either score. Joint cosine
  showed the predicted shrinkage on mild data (P4).
\item \textbf{What linear encoders cannot do} is let the band width depend on the past. On
  heavy tails every linear model is too wide on calm cases and too narrow on wild ones.
\end{itemize}

\subsection{Goal 2: nonlinear encoders}

\begin{itemize}
\item \textbf{The best contrastive models get within about 1\% of perfect on every series},
  with calibrated bands: \Ltwo{} joint and normalised cosine (Table~\ref{tab:synthetic}).
\item \textbf{Cosine's spread flaw mostly disappears.} On the Gaussian series cosine went
  from +5.6\% to +1.4--2.7\%, and its 90\% bands on the hardest cases from 98\% to 89\%
  coverage. So P1 is a statement about \textbf{linear encoders}, not about the dot product.
\item \textbf{Band widths now follow the past} on heavy tails. A regression model with a
  fixed noise band cannot do this.
\item \textbf{The encoder learns what matters without being told.} On the bent series the
  learned features recover the hidden ``un-bent'' quantity that predicts the future almost
  perfectly ($R^2 \approx 1$), and ignore the rest of the past.
\item \textbf{The shape advantage is real but conditional.} On the double-well, the best
  Gaussian MLP (same networks, past-dependent width) is at +5--6\% while \Ltwo{} is within
  1\%: a Gaussian puts probability between the two wells, where the future never goes. On
  the bent series, which is skewed but single-peaked, the Gaussian MLP ties with \Ltwo.
\item \textbf{Unbounded scores still fail on heavy tails.} Unnormalised cosine is the worst
  contrastive model there (+7\% to +23\%) and still collapses onto single futures for some
  inputs. Normalised cosine does not.
\item \textbf{Normalisation needs a nonlinear encoder.} Linear-then-normalise keeps only the
  \emph{direction} of the past and loses its \emph{size}. The size (the current level) is
  most of the information in these series, so linear normalised cosine is poor (+12--23\%).
\end{itemize}

\begin{table}[h]
\centering
\small
\begin{tabular}{@{}llccccc@{}}
\toprule
domain & dial & cos 2-sided & norm.\ cos & \Ltwo{} joint & MSE + noise & Gauss NLL \\
\midrule
Gaussian (GP)  & ---          & $5.6 \to 1.4$ & $13.1 \to 0.5$ & $0.0 \to 0.1$ & $0.1 \to 0.6$ & $0.2 \to 0.4$ \\
Heavy tails    & df $=3$      & ${\sim}10\times \to 23.3$ & $23.0 \to \mathbf{0.9}$ & $4.5 \to \mathbf{0.7}$ & $8.6 \to 9.5$ & $45.9 \to 22.2^{*}$ \\
Bent GP        & full         & $40.0 \to 3.9$ & $12.2 \to 0.9$ & $1.6 \to 0.7$ & $12.7 \to 14.7$ & $27.5 \to 0.7$ \\
Double-well    & $\lambda=1.5$ & $2.6 \to 0.3$ & $0.7 \to 0.3$ & $0.3 \to 0.0$ & $14.1 \to 7.6$ & $13.5 \to 5.2$ \\
\bottomrule
\end{tabular}
\caption{Synthetic series: energy score above the perfect forecaster, in \%,
linear $\to$ MLP encoders (mean of 2 seeds; 0 = perfect). $^{*}$One seed $+0.8$, the
other $+43.6$. Full grids are in \texttt{goal1\_results.md} and \texttt{goal2\_results.md}.}
\label{tab:synthetic}
\end{table}

\subsection{Goal 3: real weather data}

\begin{itemize}
\item \textbf{The probabilistic view works with no perfect forecaster to lean on.} When
  the next 6\,h of temperature is forecast as a \emph{change from the last reading},
  linear \Ltwo{} joint is the best model in the study, and calibrated. Its energy score is
  48\% below that of ``repeat the last value'', just ahead of the Gaussian model built the
  same way (47\%) and MSE (46\%); see Table~\ref{tab:weather}.
\item \textbf{Linear cosine shows P1 at full size on real data}: right average, 90\% bands
  about $5\times$ too wide (20--23\,\textdegree C against 4\,\textdegree C for \Ltwo).
\item \textbf{Rain (spiky, mostly zero) is the real-data heavy-tail test, and it behaves as
  goals 1--2 predicted.} Unnormalised cosine collapses onto single futures. The Gaussian
  models put about 40\% of their samples on \emph{negative} rain and say ``it will rain''
  for nearly every window. \Ltwo{} and normalised cosine keep ``mostly dry, sometimes a
  burst'', with flat calibration histograms. But nothing forecasts rain well from these
  inputs: the best models beat ``random past weather'' by about 1\%.
\item \textbf{Two real-data limits that the synthetic rounds did not show.}
  \begin{itemize}
  \item \emph{Little data.} The MLPs over-fit within a few hundred to a few thousand steps
    and end up over-confident, even with early stopping. On temperature, the MLP version
    of each strong model (\Ltwo, MSE, Gaussian NLL) was worse than the linear one; only the
    cosine scores, which are poor with linear encoders, gained. One year of heavily overlapping 10-minute windows is only about 190
    days of distinct weather.
  \item \emph{The bank.} At 10 minutes ahead the contrastive forecast was worse than
    ``repeat the last value''. With 5.8k bank futures, few start at exactly today's
    temperature, so the first steps are too spread out. Two fixes confirm this is the
    bottleneck: a bigger bank ($+0.38 \to +0.45$ skill), and forecasting the change from
    the last reading ($+0.38 \to +0.48$), which lets futures from cold days serve warm days.
  \end{itemize}
\item \textbf{Modality ablation.} The target's own past is almost everything. Radiation and
  moisture help most of the rest, mostly by narrowing the bands. Dropping temperature's own
  past made the contrastive forecast \emph{over}-confident rather than more cautious (open
  question~4).
\end{itemize}

\begin{table}[h]
\centering
\small
\begin{tabular}{@{}lcccccc@{}}
\toprule
& \multicolumn{2}{c}{temperature (level)} & \multicolumn{2}{c}{temperature (change)} & \multicolumn{2}{c}{rain} \\
& \multicolumn{2}{c}{skill vs persistence} & \multicolumn{2}{c}{skill vs persistence} & \multicolumn{2}{c}{skill vs climatology} \\
\cmidrule(lr){2-3}\cmidrule(lr){4-5}\cmidrule(lr){6-7}
model & linear & MLP & linear & MLP & linear & MLP \\
\midrule
\Ltwo{} joint   & $+0.38$ & $\mathbf{+0.28}$ & $\mathbf{+0.48}$ & $\mathbf{+0.46}$ & $\mathbf{+0.01}$ & $-0.04$ \\
norm.\ cos      & $-0.30$ & $+0.07$ & $+0.27$ & $+0.38$ & $-0.00$ & $-0.06$ \\
cos 2-sided     & $-0.56$ & $+0.14$ & $+0.21$ & $+0.42$ & $-0.92$ & $-1.04$ \\
MSE + noise     & $+0.44$ & $+0.25$ & $+0.46$ & $+0.41$ & $-1.02$ & $-1.44$ \\
Gauss NLL       & $\mathbf{+0.46}$ & $+0.21$ & $+0.47$ & $+0.41$ & $-1.16$ & $-0.67$ \\
\bottomrule
\end{tabular}
\caption{Weather, test period (late October--December 2020), mean of 2 seeds. Skill is
the fractional reduction in energy score against the reference (higher is better; 0 = no
better). For rain, climatology is the meaningful reference: it already beats persistence.
Full tables are in \texttt{test\_bed/results/weather/}.}
\label{tab:weather}
\end{table}

%======================================================================
\section{Scorecard: how the paper transfers}

\begin{table}[h]
\centering
\small
\begin{tabularx}{\textwidth}{@{}>{\raggedright\arraybackslash}p{3.6cm}LLLL@{}}
\toprule
prediction & Gaussian, linear (paper's setting) & non-Gaussian, linear & nonlinear encoders & real data \\
\midrule
Model = reweighted real futures is a usable forecaster
  & \yes{} matches closed form & \yes & \yes{} within ${\sim}1\%$ of perfect & \yes{} best model in study \\
P2: \Ltwo{} gets average \emph{and} spread right
  & \yes{} exactly & \yes{} 0--5\% of perfect & \yes & \yes{} calibrated (linear); MLPs over-confident \\
P1: dot product gets spread wrong
  & \yes{} bands too wide & \yes{} worse; breaks on heavy tails & \no{} mostly fixed & \yes{} linear ($5\times$ too wide); collapses on rain \\
P3: one-way vs two-way doesn't matter
  & \yes & \yes{} (for \Ltwo{} too) & \yes{} noisy on heavy tails & \yes{} linear; MLPs vary \\
P4: joint shrinks, worse conditional forecaster
  & \yes{} for cosine & \yes{} on mild data; \emph{less} broken on heavy tails & \no{} joint best or tied & \no{} joint best \\
P5: retrieval capped by information
  & \yes{} \Ltwo{} on the ceiling & --- & --- & --- \\
\bottomrule
\end{tabularx}
\caption{Each prediction across the four settings.}
\label{tab:scorecard}
\end{table}

Reading across the rows:
\begin{itemize}
\item \textbf{What transfers cleanly} is the framework (a contrastive model \emph{is} a
  conditional distribution) and the \Ltwo{} result. They survive every relaxation we tried.
  The framework's main practical benefit, a forecast whose shape comes from real data, is
  invisible in the paper's Gaussian setting. It becomes the main argument for the method on
  real, non-Gaussian data.
\item \textbf{What transfers only in its own setting} is the cosine--\Ltwo{} contrast (P1)
  and the joint-loss penalty (P4). Both are properties of what a \emph{linear} encoder can
  represent. A flexible encoder can compensate for the dot product's missing term, and the
  gap closes.
\item \textbf{What the paper does not address, and which turned out to matter most} off its
  setting: whether the score is bounded (heavy tails); the bank of futures the forecast is
  built from; finite data and over-confidence; and the representation of the target (level
  vs change).
\end{itemize}

%======================================================================
\section{Rules of thumb: what to expect where}

From the three rounds together, these are the predictions we would now make for a new
problem.
\begin{itemize}
\item \textbf{Linear encoders:} use \Ltwo. Expect the dot product to give the right average
  with bands far too wide. Do not normalise: it throws away the level.
\item \textbf{Flexible encoders, well-behaved data:} the choice of score barely matters.
  \Ltwo, joint cosine and normalised cosine all end up close to the best possible.
\item \textbf{Heavy-tailed or spiky targets} (rain, bursts, outliers): use a bounded score,
  i.e.\ \Ltwo{} or \emph{normalised} cosine. Expect the unnormalised dot product to collapse
  onto a few extreme futures, and Gaussian baselines to put mass on impossible values.
\item \textbf{Near-Gaussian targets} (6-hour temperature): expect contrastive and Gaussian
  forecasters to be close. The shape advantage only pays off when the future has shape:
  two outcomes, a floor at zero, strong skew.
\item \textbf{Sharp, short-range forecasts:} expect the bank to be the bottleneck. Use as
  many futures as possible, and a target that transfers between situations (changes rather
  than levels).
\item \textbf{Small or highly overlapping data:} expect MLP encoders to become over-confident
  quickly. Early stopping on a held-out probabilistic score is necessary, and on our weather
  data it was still not enough.
\end{itemize}

%======================================================================
\section{What matters and what doesn't}

\paragraph{Matters.}
\begin{itemize}
\item \textbf{The score family, when the encoder is linear.} \Ltwo{} vs dot product is the
  difference between calibrated bands and bands $5\times$ too wide.
\item \textbf{Whether the score is bounded, when the data has heavy tails.} This decides
  collapse vs no collapse, whatever the encoder.
\item \textbf{The bank:} its size, and whether its futures can stand in for the case at
  hand. Worth 7--18 points of skill on real temperature.
\item \textbf{Encoder flexibility relative to data size.} Flexibility helped a lot on
  plentiful synthetic data and hurt on one year of weather.
\item \textbf{How you evaluate.} The old metrics under-reported \Ltwo{} about $5\times$ on
  retrieval, invented a ``theory floor'' for cosine, and could not see uncertainty at all.
\end{itemize}

\paragraph{Doesn't matter (much).}
\begin{itemize}
\item \textbf{One-way vs two-way training.} It made no consistent difference, as the paper
  predicts for the dot product. Empirically the same holds for \Ltwo{} (MLPs on real data
  vary more between the two, with no consistent winner).
\item \textbf{Cosine normalisation, in most settings.} On well-behaved data with MLP
  encoders it neither helps nor harms. It matters in exactly two places, in opposite
  directions: it \emph{hurts} with linear encoders (it loses the level), and it is
  \emph{needed} for the dot product on heavy-tailed data.
\item \textbf{Whether the bank overlaps the training data.} We expected overlap to make MLPs
  look over-confident. On weather it did not; size mattered more.
\item \textbf{Batch size, for cosine's error against the theory.} There is no trend from 64
  to 1024 negatives.
\end{itemize}

%======================================================================
\section{The normalisation story: where it stands}

\paragraph{The observation (solid).} Across goals 1--3, every catastrophic failure is an
unbounded score on heavy-tailed or spiky data:
\begin{itemize}
\item linear cosine on heavy tails (about $10\times$ worse than perfect);
\item MLP cosine on heavy tails (+7\% to +23\%);
\item unnormalised cosine on rain: linear (energy 1.4--1.9$\times$ climatology's) and
  one- and two-way MLP (2.0--2.4$\times$).
\end{itemize}
The bounded scores (\Ltwo{} and normalised cosine) never fail like that. One caveat: some
MLP \Ltwo{} forecasts on rain also lean on only a couple of bank futures, but without the
loss of skill.

\paragraph{The mechanism (plausible, not yet measured).} A dot-product score grows with the
size of the future. Reweighting by $\exp(\text{score})$ therefore favours the biggest
futures without limit, and on heavy-tailed data a few extreme futures take all the weight.
CLIP also normalises for practical reasons (stable training, a fixed temperature range), so
the observation alone does not prove this is \emph{the} reason.

\paragraph{A direct check (cheap).} For unnormalised cosine on rain, measure the length of
each future's embedding against the size of that future, and how concentrated the forecast
weights are. If the long embeddings belong to the spiky futures and those futures take the
weight, the mechanism is confirmed.

\paragraph{A clean theoretical argument (likely exists).} The model distribution is
\[
  \nu(v \mid u) \;=\; \frac{e^{s(u,v)}\,\mu_v(v)}{Z(u)},
  \qquad
  Z(u) \;=\; \int e^{s(u,v')}\,\mu_v(dv'),
\]
i.e.\ $\exp(\text{score})$ times the distribution of futures $\mu_v$, rescaled to sum to 1.
\begin{itemize}
\item For a score that is linear in the future, $s(u,v) = \langle a(u), v\rangle$ (the
  dot product with a linear future encoder), $Z(u)$ is the moment generating function of
  $\mu_v$. It is finite only if the futures have light, exponential-type tails. For a
  Student-$t$ it is infinite for every $a(u) \neq 0$, so the ``model distribution'' the loss
  is supposed to be learning does not exist. A finite batch or bank hides this by collapsing
  onto its largest members, which is exactly what we observe.
\item A score that is \textbf{bounded above}, $s \le c$, gives $Z(u) \le e^{c}$ for
  \emph{any} distribution of futures. \Ltwo{} is at most 0 and normalised cosine at most
  its scale, so both always define a valid model.
\end{itemize}
So bounded scores are the ones for which the paper's probabilistic objective is valid on
arbitrary data. That is a stronger statement than ``CLIP normalises for a reason'', it comes
straight out of the paper's framework, and it could be written up properly.

%======================================================================
\section{A research direction: the bank (reference set)}

Forecasting by reweighting a bank of real futures is \textbf{importance sampling}. A known
result~\cite{cd} says the number of samples needed grows roughly like
$\exp(\mathrm{KL})$, where KL measures how far the target distribution (here, the forecast)
is from the distribution the samples come from (here, the bank). That explains the
short-range failure in goal~3:
\begin{itemize}
\item At 10 minutes ahead the true forecast is very sharp, so it is far from the bank's
  broad spread. KL is large, and 5.8k futures are not enough.
\item A bigger bank raises the sample count.
\item The change-from-last-value target shrinks the gap, because a bank of \emph{changes}
  is much closer to any particular forecast than a bank of \emph{levels}.
\end{itemize}
A theory version would split forecast error into \textbf{score error} (did training find
the right weights?) plus \textbf{reference-set error}: bank size, a shift between the bank's
period and the test period, and how the future is represented. It would also say which
representations make a bank transferable. This lies outside the paper, which treats the
bank as the true data distribution, and it is the limit that bit hardest on real data.

%======================================================================
\section{Open questions}

\begin{enumerate}
\item \textbf{Why does cosine work with MLP encoders?} The guess from goal~2 is that MLP
  bias terms let the dot product contain a piece that depends on the future alone. That is
  the piece the linear dot product is missing (the reason for P1) and that \Ltwo{} has
  built in. The paper's own footnote on affine encoders points the same way.
  \emph{Cheap test:} add an explicit future-only term to linear cosine and see whether the
  spread is fixed.
\item \textbf{Over-sharpening.} Longer training made MLP forecasts over-confident: in
  goal~2, 30k steps were worse than 15k; in goal~3, MLPs over-fit within a few hundred to a
  few thousand steps. Is the learned temperature (the logit scale) the culprit?
  \emph{Cheap test:} after training, tune only the temperature on validation. If that fixes
  calibration, the objective has a one-parameter sharpness problem, a cleaner story than
  ``use early stopping''.
\item \textbf{Why is joint consistently best?} \Ltwo{} joint was the best or tied for best
  in goals 2 and 3. On rain, joint cosine was the only unnormalised MLP cosine that did not
  collapse, and in goal~1 it was much less broken on heavy tails. The paper separates joint
  from conditional losses, but predicts the opposite ranking for conditional forecasting
  (P4). One guess: the joint softmax's shrinkage damps the pull towards extreme futures.
  Untested.
\item \textbf{Over-confidence when the input says little.} With temperature's own past
  removed, the contrastive forecast became over-confident instead of falling back towards
  ``random past weather''. Whether that is an estimation artefact or structural matters for
  real cross-modal use, where one modality often says little about the other.
\item \textbf{Does ``the encoder learns the useful quantity'' (goal~2) hold on real data?}
  We never checked what the weather encoders learned.
\end{enumerate}

%======================================================================
\section{Caveats}

\begin{itemize}
\item Synthetic rounds: 2 seeds and 500 test cases per setting, so differences under about
  1\% are noise. The fixed synthetic field setup has observation noise larger than the
  signal's own variation, which mutes the bent-series effects.
\item Real data: one year, one station, one test period (late October--December, colder
  than most of training), 2 seeds.
\item The MLPs were barely tuned: one architecture, no weight decay or dropout, and only the
  learning rate compared. A regularised MLP may well beat the linear models on weather; that
  is the obvious next experiment.
\item The normalisation mechanism and every open question are hypotheses, not results.
\end{itemize}

\begin{thebibliography}{9}
\bibitem{bst} R.~Baptista, A.~Stuart and S.~Tran.
  \emph{A Mathematical Perspective on Contrastive Learning}. arXiv:2505.24134, 2025.
\bibitem{cd} S.~Chatterjee and P.~Diaconis.
  The sample size required in importance sampling.
  \emph{Annals of Applied Probability}, 28(2):1099--1135, 2018.
\bibitem{gneiting} T.~Gneiting and A.~E.~Raftery.
  Strictly proper scoring rules, prediction, and estimation.
  \emph{Journal of the American Statistical Association}, 102(477):359--378, 2007.
\end{thebibliography}

\end{document}
